% =============================================================================
%  Corretude e desempenho
%  Cobre os subprojetos testing, undefined_behavior e performance.
% =============================================================================

% =============================================================================
\section{Corretude e desempenho}
\label{sec:qualidade}
% =============================================================================

Um programa em C++ pode errar de três maneiras distintas, e elas exigem
ferramentas distintas. Pode dar a \textbf{resposta errada} --- e o instrumento é
o teste automatizado. Pode não ter resposta \emph{nenhuma} definida pelo padrão
--- é o \gls{ub}, e o instrumento é o \emph{sanitizer}, o aviso do compilador e,
sobretudo, o estilo de código que não dá chance ao erro. E pode dar a resposta
certa \textbf{tarde demais} --- e o instrumento é a medição. Este capítulo
percorre os três: \projeto{testing} (\cref{sec:testes}),
\projeto{undefined\_behavior} (\cref{sec:ub}) e \projeto{performance}
(\cref{sec:desempenho}).

O que amarra os três é um método, não um assunto. Eles \textbf{medem} em vez de
afirmar: a matriz de detecção dos \emph{sanitizers} não foi copiada de
documentação, saiu de rodar código com UB de propósito neste \emph{toolchain};
os fatores de $35\times$ e $11{,}7\times$ saíram de um cronômetro. E, em vários
pontos, \textbf{o número contrariou a expectativa de quem escreveu o texto}.
Esses casos estão todos registrados aqui, porque um resultado que desmente a
hipótese ensina método; um que a confirma só ensina o assunto.

\begin{porque}[por que corretude e desempenho estão no mesmo capítulo]
Porque as três disciplinas se cobrem mutuamente e nenhuma fecha sozinha. O teste
encontra a lógica errada e é cego para o \gls{ub} que ainda não se manifestou. O
\emph{sanitizer} encontra o UB, mas só no caminho que \emph{de fato} executou ---
por isso depende dos testes para ter o que executar. A cobertura diz qual caminho
\textbf{não} rodou, que é exatamente o ponto cego do \emph{sanitizer}. E a
medição de desempenho só faz sentido sobre um programa já correto: acelerar
código errado é produzir a resposta errada mais depressa.
\end{porque}

% =============================================================================
\subsection{Testes automatizados}
\label{sec:testes}
% =============================================================================

O subprojeto \projeto{testing} é o único do repositório com \textbf{dois tipos de
alvo}: um executável didático, que narra o assunto, e arquivos
\texttt{test\_*.cpp} que são testes de verdade, rodados por \texttt{meson test}.
Estado atual: 6 alvos, 41 casos, todos passando.

\begin{porque}[que problema o teste automatizado resolve]
Não é ``garantir que o programa funciona'' --- nenhuma suíte prova isso. É
\textbf{tornar barata a mudança}. Sem testes, cada alteração exige reler
mentalmente tudo o que ela pode ter quebrado, e o custo dessa releitura cresce
com o sistema até o ponto em que ninguém mais mexe no que funciona. Com testes, a
pergunta ``eu quebrei alguma coisa?'' vira um comando de dez segundos. A suíte
não é prova de corretude; é \textbf{rede de segurança para a refatoração}.
\end{porque}

A pergunta útil, portanto, não é ``devo testar?''. É \textbf{``que teste vale a
pena escrever?''} --- porque uma suíte grande de testes ruins é \emph{pior} que
nenhuma: custa manutenção, falha por motivos errados e ensina a equipe a ignorar
a luz vermelha.

\subsubsection{A pirâmide, e por que ela é uma pirâmide}

Três definições primeiro, porque os termos são usados de forma descuidada. Um
\textbf{teste unitário} exercita uma unidade --- tipicamente uma classe ou função
--- isolada de tudo o que ela depende: sem banco, sem rede, sem arquivo, sem
relógio; roda em microssegundos. Um \textbf{teste de integração} exercita duas ou
mais peças juntas, com as dependências reais. Um \textbf{teste de ponta a ponta}
(E2E) exercita o sistema inteiro pela mesma porta que o usuário usa.

\begin{figure}[htbp]
\centering
\begin{tikzpicture}[node distance=2.5mm]
  \node[caixarisco, minimum width=32mm] (e2e)
      {E2E / sistema\\\scriptsize poucos --- lentos, frágeis, caros};
  \node[caixa, minimum width=54mm, below=of e2e] (int)
      {integração\\\scriptsize alguns --- banco, rede, IO};
  \node[caixaok, minimum width=78mm, below=of int] (uni)
      {UNITÁRIOS\\\scriptsize muitos --- rápidos, estáveis, baratos};
\end{tikzpicture}
\caption{A pirâmide de testes. A forma é consequência do custo por teste em três
         eixos: \textbf{tempo} (microssegundos contra minutos),
         \textbf{estabilidade} (mais peças, mais motivos de falhar sem haver bug)
         e \textbf{diagnóstico} (um unitário aponta a linha; um E2E aponta
         ``algo, em algum lugar'').}
\label{fig:piramide}
\end{figure}

\begin{armadilha}[a casquinha de sorvete]
A pirâmide invertida acontece quando a equipe testa só pela interface do usuário,
porque é ali que ``dá para ver''. O sintoma é conhecido: a suíte demora 40
minutos e falha uma vez em cinco sem motivo. A partir daí ninguém confia no
vermelho, e a suíte inteira vira decoração de \emph{pipeline}.
\end{armadilha}

\subsubsection{F.I.R.S.T. e o que testar primeiro}

\begin{table}[htbp]
\centering
\caption{F.I.R.S.T.: o que cada letra exige, e o que a viola.}
\label{tab:first}
\footnotesize
\begin{tabular}{@{}lll@{}}
\toprule
\textbf{Letra} & \textbf{Exigência} & \textbf{Inimigo típico} \\
\midrule
\textbf{F}ast            & a suíte inteira roda em segundos & IO, \emph{sleep}, banco real \\
\textbf{I}solated        & independe dos outros e da \textbf{ordem} & estado global, \emph{singleton} \\
\textbf{R}epeatable      & mesmo resultado em qualquer máquina & relógio, rede, sorteio sem semente \\
\textbf{S}elf-validating & passa ou falha, sem ninguém ler a saída & teste que pede conferência humana \\
\textbf{T}imely          & escrito junto com o código & ``a gente testa depois'' \\
\bottomrule
\end{tabular}
\end{table}

A lista de inimigos do \textbf{R} é maior do que parece: além do relógio e da
rede, entram a ordem de iteração de um \classe{std::unordered\_map} --- que não é
especificada e muda entre versões da biblioteca ---, caminhos absolutos e fuso
horário. E o \textbf{T} carrega a lição mais valiosa do subprojeto:
\textbf{código difícil de testar é código mal desenhado}, e o teste escrito cedo
devolve esse recado a tempo de agir.

Em ordem de retorno sobre o esforço, teste: (1)~\textbf{as fronteiras} --- a
regra de desconto do sistema sob teste tem faixas em 10, 50 e 100 peças, e os
casos que importam são 9/10, 49/50, 99/100, porque é ali que mora o \texttt{>}
que deveria ser \texttt{>=}; (2)~\textbf{o caminho de erro}, que é o menos
exercitado manualmente e por isso o mais quebrado; (3)~\textbf{o bug que você
acabou de consertar} --- um teste por bug corrigido é a regra de melhor
custo-benefício que existe, porque ele nunca volta; (4)~\textbf{a regra de
negócio}, aquilo que, errado, cobra do cliente o valor errado; (5)~o caminho
feliz, por último. E não teste \emph{getters} triviais, a biblioteca padrão,
detalhes privados de implementação nem código gerado.

\begin{armadilha}[os cinco testes que dão prejuízo]
\textbf{(a) O frágil} --- falha em toda refatoração \emph{correta}; sintoma:
cheio de \texttt{EXPECT\_CALL}. Testa \emph{como}, não \emph{o quê}.
\textbf{(b) O instável} (\emph{flaky}) --- passa às vezes. É o pior de todos,
porque destrói a confiança na suíte inteira: depois de dois falsos alarmes, todo
mundo reroda em vez de investigar. Deve ser \textbf{consertado ou removido},
nunca deixado piscando.
\textbf{(c) O tautológico} --- repete a implementação:
\texttt{EXPECT\_EQ(c.total(), c.subtotal() * (1 - c.desconto()))} passa mesmo com
a fórmula errada; o valor esperado precisa vir \textbf{de fora}, como literal.
\textbf{(d) O sem asserção} --- só verifica que não explodiu.
\textbf{(e) O gigante} --- quando falha a asserção nº 4, você não descobre se as
outras 26 falhariam.
\end{armadilha}

\subsubsection{O sistema sob teste}

Tudo o que vem a seguir gira em torno de uma classe pequena, em
\texttt{src/testing/carrinho.hpp}, escolhida por ter de propósito os cinco
problemas que testes de verdade encontram: uma regra de negócio com faixas, uma
pré-condição que aborta, uma exceção, uma dependência externa e aritmética de
ponto flutuante.

\begin{lstlisting}[style=bt, language=C++,
    caption={O sistema sob teste, resumo de \texttt{src/testing/carrinho.hpp}.}]
class RepositorioDePrecos {                     // a DEPENDÊNCIA, como interface
    virtual std::optional<double> preco_de(const std::string& sku) const = 0;
};

class Carrinho {
    // A dependência entra pelo CONSTRUTOR. É esta linha que torna a classe
    // testável em microssegundos, sem banco de dados.
    explicit Carrinho(const RepositorioDePrecos& precos) : precos_{precos} {}

    void adicionar(const std::string& sku, int quantidade) {
        // PRÉ-CONDIÇÃO: quantidade <= 0 é BUG DO CHAMADOR, não erro de entrada.
        assert(quantidade > 0 && "quantidade precisa ser positiva");
        const auto preco = precos_.preco_de(sku);
        if (!preco) throw ProdutoDesconhecido{sku};   // ← erro de ENTRADA
    }
    double percentual_de_desconto() const noexcept;  // faixas: 10, 50, 100
};
\end{lstlisting}

A distinção entre as duas espécies de coisa ruim --- erro de entrada, que vira
exceção, e bug do chamador, que vira aborto --- é o assunto do \cref{sec:errors},
e aqui tem consequência direta sobre \emph{como} se testa cada uma.

\subsubsection{Anatomia de um teste}

Um caso tem a forma \texttt{TEST(Suite, Caso) \{ ... \}}, e o par
\texttt{(Suite, Caso)} precisa ser único no binário inteiro. Dentro do corpo, a
organização que funciona tem três blocos separados por uma linha em branco:
\textbf{Arrange} (preparar), \textbf{Act} (a ação testada) e \textbf{Assert}.

\begin{lstlisting}[style=bt, language=C++,
    caption={Arrange / Act / Assert, em \texttt{src/testing/test\_basico.cpp}.}]
TEST(Carrinho, MesmoSkuDuasVezesSomaQuantidade) {
    PrecosFixos    precos;                  // ARRANGE
    loja::Carrinho c{precos};

    c.adicionar("caneta", 2);               // ACT
    c.adicionar("caneta", 5);

    ASSERT_EQ(c.tamanho(), 1u) << "deveria continuar sendo UM item";
    EXPECT_EQ(c.itens()[0].quantidade, 7);  // ASSERT
}
\end{lstlisting}

Se o \emph{Act} tem mais de duas ou três linhas, ou se há \emph{dois} blocos
act-assert no mesmo caso, o teste está testando duas coisas --- separe.

\begin{regra}[num teste, a repetição é preferível à abstração]
É o único código em que isso vale. \texttt{EXPECT\_DOUBLE\_EQ(total, 77.9)} com o
número literal é \textbf{melhor} que uma função \texttt{calcula\_esperado()} ---
que repetiria a implementação e passaria a esconder exatamente o bug que o teste
deveria pegar. O valor esperado precisa entrar \emph{de fora}, calculado por
outro meio que não o código sob teste.
\end{regra}

O nome do caso é lido por alguém que não o escreveu, no relatório de uma falha. A
fórmula que funciona é \texttt{<situação><ação><resultado esperado>}:
\texttt{TEST(Desconto, ExatamenteDezPecasDaCincoPorCento)} contra
\texttt{TEST(CarrinhoTest, Test1)} ou \texttt{FuncionaCorretamente}. Os nomes
viram \textbf{identificadores C++}: nada de espaço, acento ou hífen --- e nada de
sublinhado, que o próprio GoogleTest desaconselha, porque ele usa \texttt{\_}
internamente para colar suíte e caso.

\subsubsection{\texttt{EXPECT} contra \texttt{ASSERT}, e o catálogo}

A diferença é de uma palavra e muda o fluxo: \texttt{EXPECT\_*} registra a falha e
\textbf{continua}; \texttt{ASSERT\_*} registra e faz \texttt{return}.

\begin{regra}[\texttt{EXPECT} por padrão, \texttt{ASSERT} quando o resto não faz sentido]
\texttt{EXPECT} dá mais informação por execução. Use \texttt{ASSERT} quando
continuar seria inútil ou perigoso --- o caso clássico é verificar o tamanho
antes de indexar: \texttt{ASSERT\_EQ(c.tamanho(), 1u);} e só então
\texttt{EXPECT\_EQ(c.itens()[0].sku, "caneta");}. Sem o \texttt{ASSERT}, a linha
seguinte indexaria um vetor vazio, o que é \gls{ub}.
\end{regra}

\begin{armadilha}[\texttt{ASSERT\_*} faz \texttt{return}, e isso tem três consequências]
(1)~Não funciona em função que devolva valor --- o \texttt{return} é sem
argumento. (2)~Dentro de uma \textbf{função auxiliar} ele retorna da
\emph{auxiliar}, e o teste continua como se nada tivesse acontecido; para
propagar, envolva a chamada em
\texttt{ASSERT\_NO\_FATAL\_FAILURE(minha\_verificacao())}. (3)~Não pode ser usado
em construtor nem em destrutor --- é para isso que existem \texttt{SetUp()} e
\texttt{TearDown()}.
\end{armadilha}

\begin{table}[htbp]
\centering
\caption{As famílias de asserções. Todas têm a variante \texttt{ASSERT\_}.}
\label{tab:assercoes}
\footnotesize
\begin{tabular}{@{}lll@{}}
\toprule
\textbf{Família} & \textbf{Macros} & \textbf{Observação} \\
\midrule
booleanas   & \texttt{EXPECT\_TRUE} / \texttt{FALSE} & imprime só ``false'' na falha \\
comparação  & \texttt{EXPECT\_EQ NE LT LE GT GE}     & imprimem \textbf{os dois} valores \\
strings C   & \texttt{EXPECT\_STREQ STRNE STRCASEEQ} & para \texttt{std::string}, use \texttt{EXPECT\_EQ} \\
ponto flut. & \texttt{EXPECT\_FLOAT\_EQ DOUBLE\_EQ NEAR} & 4 ULPs, ou tolerância explícita \\
exceções    & \texttt{EXPECT\_THROW NO\_THROW ANY\_THROW} & o tipo é o segundo argumento \\
morte       & \texttt{EXPECT\_DEATH DEBUG\_DEATH}    & roda em processo filho \\
controle    & \texttt{FAIL() ADD\_FAILURE() SUCCEED() GTEST\_SKIP()} & --- \\
\bottomrule
\end{tabular}
\end{table}

A razão de preferir \texttt{EXPECT\_EQ(a, b)} a \texttt{EXPECT\_TRUE(a == b)}
está na terceira coluna: o primeiro imprime os dois valores quando falha, o
segundo imprime ``false''. A mesma lógica leva aos \emph{matchers} do gmock,
usáveis em testes comuns via \texttt{EXPECT\_THAT} --- um \texttt{ElementsAre} que
falha imprime os dois vetores inteiros e marca o índice divergente:

\begin{lstlisting}[style=bt, language=C++,
    caption={\emph{Matchers}: valem sobretudo pela mensagem de falha.}]
EXPECT_THAT(v, ElementsAre(1, 2, 3));      EXPECT_THAT(v, SizeIs(Ge(3)));
EXPECT_THAT(v, UnorderedElementsAre(3, 1, 2));  EXPECT_THAT(s, HasSubstr("erro"));
EXPECT_THAT(x, AllOf(Gt(0), Lt(100)));     EXPECT_THAT(o, Optional(Eq(42)));
\end{lstlisting}

\subsubsection{Ponto flutuante}

Um \texttt{double} representa números em base 2, e \texttt{0.1} não tem
representação exata em base 2 --- como $1/3$ não tem representação decimal
finita. O subprojeto imprime a consequência:

\begin{lstlisting}[style=bt, language=bash,
    caption={Saída de \texttt{src/testing/}: por que \texttt{==} não serve.}]
0.1 + 0.2 = 0.30000000000000004441
0.3       = 0.29999999999999998890
diferença = 5.551e-17
\end{lstlisting}

A unidade natural de tolerância é o \textbf{ULP} (\emph{unit in the last place}):
a distância entre dois \texttt{double} consecutivos representáveis em torno
daquele valor. Ela é minúscula perto de zero e grande perto de $10^{15}$, que é
exatamente o que se quer --- por isso \texttt{EXPECT\_DOUBLE\_EQ}, que aceita
4~ULPs, \emph{escala com a magnitude}, enquanto \texttt{EXPECT\_NEAR(a, b, 1e-9)}
usa tolerância absoluta.

\begin{armadilha}[nenhuma comparação de \emph{float} funciona bem em torno de zero]
4~ULPs de zero é um número absurdamente pequeno --- praticamente uma exigência de
igualdade exata. Perto de zero, use \texttt{EXPECT\_NEAR} com uma tolerância que
faça sentido para o \emph{domínio}. E \texttt{NaN} não é igual a nada, nem a si
mesmo: verifique com \texttt{EXPECT\_TRUE(std::isnan(x))}.
\end{armadilha}

\subsubsection{Fixtures e o ciclo de vida}

Uma \textbf{fixture} é uma classe que carrega a preparação comum a vários casos.
Deriva de \texttt{::testing::Test}, e os casos passam a usar \texttt{TEST\_F},
com o nome da \emph{classe} como primeiro argumento. A regra que mais surpreende
é a que garante o ``I'' de isolado: o GoogleTest constrói uma instância
\textbf{nova} da fixture para \textbf{cada} caso.

\begin{lstlisting}[style=bt, language=bash,
    caption={O ciclo de vida, resumo de \texttt{src/testing/demo\_organizacao.hpp}.}]
SetUpTestSuite()          ← UMA vez, antes de todos os casos
   construtor  →  SetUp()  →  caso 1  →  TearDown()  →  destrutor
   construtor  →  SetUp()  →  caso 2  →  TearDown()  →  destrutor
TearDownTestSuite()       ← UMA vez, depois de todos
\end{lstlisting}

Há, portanto, dois pares de ganchos, e a escolha entre eles não é estilística:
use o \textbf{construtor} por padrão --- ele permite membros \texttt{const} e sem
construtor padrão; use \texttt{SetUp()} quando a preparação puder falhar com
\texttt{ASSERT\_*} ou quando a fixture for herdada; use o \textbf{destrutor} por
padrão, e \texttt{TearDown()} quando a limpeza usar \texttt{ASSERT\_*} ou puder
lançar. O arquivo \texttt{src/testing/test\_fixture.cpp} tem um caso cuja única
função é \emph{provar} o isolamento --- ele falharia se a fixture fosse
compartilhada:

\begin{lstlisting}[style=bt, language=C++,
    caption={O teste que prova que a fixture é reconstruída.}]
TEST_F(CarrinhoTest, SegundoTesteVeCarrinhoLimpo) {
    EXPECT_TRUE(carrinho.vazio())
        << "se este teste falhar, a fixture está sendo compartilhada";
}
\end{lstlisting}

\begin{armadilha}[\texttt{SetUpTestSuite} compartilha estado entre casos]
Os ganchos de \emph{suíte} rodam uma única vez e são estáticos, portanto o que
eles preparam é comum a todos os casos. É a porta de entrada para o teste que só
falha numa ordem específica --- e que passa quando você o roda sozinho para
investigar. Use apenas para recursos \textbf{caros e imutáveis}.
\end{armadilha}

\subsubsection{As quatro variantes, e as FRONTEIRAS}

\begin{table}[htbp]
\centering
\caption{O que cada macro de teste faz variar.}
\label{tab:variantes-teste}
\footnotesize
\begin{tabular}{@{}lll@{}}
\toprule
\textbf{Macro} & \textbf{Varia} & \textbf{Quando usar} \\
\midrule
\texttt{TEST}       & nada    & caso isolado, sem preparação \\
\texttt{TEST\_F}    & fixture & preparação repetida \\
\texttt{TEST\_P}    & \textbf{dados} & mesma lógica, muitas entradas \\
\texttt{TYPED\_TEST}& \textbf{tipos} & \emph{template}, ou $N$ implementações da mesma interface \\
\bottomrule
\end{tabular}
\end{table}

\texttt{TEST\_P} é o que mais rende, e a razão não é economia de linhas. Um laço
\texttt{for} dentro de um \texttt{TEST\_F} produz \textbf{um} teste: ele para na
primeira falha e a mensagem diz ``algo, em algum caso''. Um
\texttt{INSTANTIATE\_TEST\_SUITE\_P} produz \textbf{N} testes independentes, e a
falha diz \texttt{Faixas/DescontoTest.AplicaAFaixaCorreta/exatamente\_50}.

\begin{lstlisting}[style=bt, language=C++,
    caption={\texttt{TEST\_P} sobre as fronteiras, em
             \texttt{src/testing/test\_parametrizado.cpp}.}]
struct CasoDesconto {          // um struct com NOMES, não um std::tuple:
    int         pecas;         // `caso.pecas` se lê melhor que get<0>(caso)
    double      desconto_esperado;
    const char* rotulo;        // vira parte do NOME do teste
};
class DescontoTest : public ::testing::TestWithParam<CasoDesconto> { /* ... */ };

TEST_P(DescontoTest, AplicaAFaixaCorreta) {
    const auto& caso = GetParam();
    loja::Carrinho c{precos};
    c.adicionar("caneta", caso.pecas);
    EXPECT_DOUBLE_EQ(c.percentual_de_desconto(), caso.desconto_esperado)
        << "com " << caso.pecas << " peças";
}

// As FRONTEIRAS são o que importa: 9/10, 49/50, 99/100. Um teste que só
// verifica 5, 30 e 200 passaria com um `>` no lugar de um `>=`.
INSTANTIATE_TEST_SUITE_P(
    Faixas, DescontoTest,
    ::testing::Values(CasoDesconto{  9, 0.00, "logo_abaixo_de_10"},
                      CasoDesconto{ 10, 0.05, "exatamente_10"},
                      CasoDesconto{ 49, 0.05, "logo_abaixo_de_50"},
                      CasoDesconto{ 50, 0.10, "exatamente_50"},
                      CasoDesconto{100, 0.15, "exatamente_100"}),
    [](const ::testing::TestParamInfo<CasoDesconto>& info) {
        return std::string{info.param.rotulo};      // o gerador de nomes
    });
\end{lstlisting}

Os geradores prontos são \texttt{Values(...)}, \texttt{Range(0, 10)} (com passo
opcional), \texttt{Bool()}, \texttt{ValuesIn(container)} --- a partir de um vetor
montado em execução --- e \texttt{Combine(a, b)}, o \textbf{produto cartesiano},
que explode rápido.

\begin{armadilha}[sem o gerador de nomes, ``o /3 falhou'' não diz nada]
O último argumento de \texttt{INSTANTIATE\_TEST\_SUITE\_P} é opcional e, na
prática, quase obrigatório: sem ele os casos se chamam \texttt{.../0},
\texttt{.../1}, e o relatório do CI vira um enigma. O \emph{lambda} que devolve o
rótulo custa três linhas.
\end{armadilha}

Os \textbf{testes tipados} resolvem o problema simétrico: a mesma lógica sobre
vários \emph{tipos} --- é como se testa um \emph{template}, ou $N$ implementações
da mesma interface.

\begin{lstlisting}[style=bt, language=C++,
    caption={\texttt{TYPED\_TEST}: um corpo, vários tipos.}]
template <typename T> class ContainerTest : public ::testing::Test {
protected: T container;
};
using TiposTestados = ::testing::Types<std::vector<int>, std::deque<int>>;
TYPED_TEST_SUITE(ContainerTest, TiposTestados);

TYPED_TEST(ContainerTest, ComecaVazioEDepoisNao) {
    EXPECT_TRUE(this->container.empty());     // ← o `this->` é obrigatório
    this->container.push_back(42);
    EXPECT_EQ(this->container.size(), 1u);
}
\end{lstlisting}

\begin{armadilha}[em \texttt{TYPED\_TEST}, os membros da fixture exigem \texttt{this->}]
A fixture é um \emph{template}, e \texttt{container} é um \textbf{nome dependente}
--- depende do parâmetro \texttt{T}. Pela regra de busca de nomes do C++, o
compilador não procura em classes-base dependentes na primeira fase da compilação
do \emph{template} (\cref{sec:templates}). Sem o \texttt{this->}, ele nem olha em
\texttt{ContainerTest<T>} e reporta ``não declarado neste escopo''.
\end{armadilha}

\begin{lstlisting}[style=bt, language=bash,
    caption={As opções de linha de comando que mais rendem.}]
--gtest_filter=*Desconto*:-*Lento*   # incluir e excluir por padrão de nome
--gtest_repeat=200 --gtest_shuffle   # a combinação que caça flakiness
--gtest_break_on_failure             # parar no depurador
--gtest_output=xml:res.xml           # para o CI
meson test -C build testing_mock -v  # um alvo, com a saída na tela
meson test -C build --repeat 50 --gdb --wrapper=valgrind
\end{lstlisting}

Um teste temporariamente quebrado se marca renomeando-o para
\texttt{DISABLED\_...}: o GoogleTest o pula e \textbf{avisa} quantos estão
desabilitados --- melhor que comentar o teste e esquecer que ele existe.

\begin{armadilha}[dois erros de infraestrutura que deixam o CI verde à toa]
\textbf{(1)} O Conan expõe quatro pacotes separados: \texttt{gtest} (asserções e
o \emph{runner}, \textbf{sem} \texttt{main}), \texttt{gtest\_main},
\texttt{gmock} e \texttt{gmock\_main}. Linkar só com \texttt{gtest} produz
\texttt{undefined reference to 'main'} --- foi o que aconteceu na primeira versão
do \texttt{meson.build} deste subprojeto.
\textbf{(2)} \texttt{RUN\_ALL\_TESTS()} precisa ter o retorno \emph{retornado}:
ignorá-lo faz o binário sair com código 0 mesmo com testes falhando, e o CI fica
verde com a suíte quebrada. O GoogleTest marca a função \texttt{[[nodiscard]]}
justamente por isso.
\end{armadilha}

\subsubsection{Os cinco dublês e a injeção de dependência}

Um \textbf{dublê de teste} é qualquer objeto que substitui uma dependência real
durante o teste. A taxonomia de Meszaros distingue cinco, e a distinção importa
porque cada um responde a uma pergunta diferente.

\begin{table}[htbp]
\centering
\caption{Os cinco dublês de teste, e o que cada um verifica.}
\label{tab:dubles}
\footnotesize
\begin{tabular}{@{}lll@{}}
\toprule
\textbf{Dublê} & \textbf{O que faz} & \textbf{Verifica} \\
\midrule
\emph{dummy} & só preenche um parâmetro; nunca é usado & nada \\
\emph{stub}  & devolve respostas prontas               & \textbf{estado} \\
\emph{spy}   & um \emph{stub} que anota o que recebeu  & os dois \\
\emph{mock}  & tem expectativas programadas            & \textbf{interação} \\
\emph{fake}  & implementação real, simplificada        & \textbf{estado} \\
\bottomrule
\end{tabular}
\end{table}

\begin{regra}[a escolha certa quase sempre é \emph{fake} ou \emph{stub}]
O \emph{mock} existe para uma pergunta específica --- ``o código chamou X?'' --- e
essa pergunta só importa quando a chamada \textbf{é} o efeito que o código
produz: enviou o e-mail, gravou o \emph{log}, publicou a mensagem. Se o efeito é
observável no estado, verifique o estado.
\end{regra}

E tudo isso só é possível por causa de uma única linha de projeto ---
\texttt{explicit Carrinho(const RepositorioDePrecos\& precos);}. Compare com o
que quase todo mundo escreve primeiro:

\begin{lstlisting}[style=bt, language=C++,
    caption={A dependência construída por dentro contra a dependência injetada.}]
class Carrinho {
    BancoDeDados db_{"localhost:5432"};   // ✗ só testável com Postgres rodando
};
class Carrinho {                          // ✓ a dependência entra de fora
public:
    explicit Carrinho(const RepositorioDePrecos& precos);
};
\end{lstlisting}

Isso é \textbf{injeção de dependência}, e não é sobre ``usar interfaces em
tudo'': é sobre inverter \emph{quem decide} a dependência concreta --- não a
classe, e sim quem a constrói. É o D de SOLID (\cref{sec:solid}). Em ordem de
preferência: parâmetro de \textbf{função} (menor acoplamento possível), de
\textbf{construtor} (quando vários métodos a usam), de \textbf{template}
(despacho estático, sem \gls{vtable} e sem alocação, \cref{sec:templates}).

\begin{armadilha}[os cinco padrões que impedem o teste]
\textbf{(a) \emph{Singleton}} --- amarra o código a uma configuração global
compartilhada entre testes, e dois testes com configurações diferentes passam a
depender da ordem. \textbf{(b) \texttt{new} dentro da classe} --- se o objeto
constrói a própria dependência, não há como substituí-la. \textbf{(c) O relógio}
--- \texttt{std::chrono::system\_clock::now()} dentro da lógica torna impossível
testar a virada do mês; injete o relógio. \textbf{(d) Funções estáticas livres
que fazem IO} --- não existe ponto de extensão. \textbf{(e) Métodos privados que
você ``precisa'' testar} --- é sintoma: ou o comportamento é alcançável pela
interface pública, ou aquilo é outra classe querendo nascer. E
\texttt{\#define private public} antes do \texttt{\#include} ``funciona'' e é
\gls{ub}: muda o \emph{layout} do objeto e viola a ODR --- a Regra da Definição
Única, que exige que todo tipo tenha \emph{a mesma} definição em todas as
unidades de tradução.
\end{armadilha}

\subsubsection{gmock: comportamento e interação}

Declarar um \emph{mock} custa uma linha por método virtual. Os parênteses em
volta dos parâmetros são obrigatórios, e um tipo com vírgula ---
\texttt{std::map<int,int>} --- precisa dos seus próprios parênteses ou de um
\emph{alias}, senão o pré-processador quebra a lista de argumentos na vírgula
errada.

\begin{lstlisting}[style=bt, language=C++,
    caption={\texttt{MOCK\_METHOD} e \texttt{EXPECT\_CALL}, em
             \texttt{src/testing/test\_mock.cpp}.}]
class MockPrecos : public loja::RepositorioDePrecos {
public:
    // MOCK_METHOD(retorno, nome, (parâmetros), (qualificadores));
    MOCK_METHOD(std::optional<double>, preco_de, (const std::string& sku),
                (const, override));
};

TEST(MockPrecos, ConsultaUmaVezPorAdicionar) {
    MockPrecos precos;
    // Esta é a asserção do teste, e ela é verificada no DESTRUTOR do mock:
    // não há EXPECT_ nenhum no corpo, e ainda assim o teste pode falhar.
    EXPECT_CALL(precos, preco_de(Eq("caneta")))
        .Times(2)
        .WillRepeatedly(Return(std::optional<double>{2.50}));

    loja::Carrinho c{precos};
    c.adicionar("caneta", 1);
    c.adicionar("caneta", 1);
}
\end{lstlisting}

São duas coisas distintas: \texttt{ON\_CALL(...).WillByDefault(...)}
\textbf{define comportamento} e não é asserção; \texttt{EXPECT\_CALL(...)}
\textbf{é a asserção}.

\begin{armadilha}[a ordem das regras de \texttt{ON\_CALL} é a inversa da intuição]
Descoberta na primeira execução de \texttt{test\_mock.cpp}:
\begin{lstlisting}[style=bt, language=C++, frame=none,
                   xleftmargin=0pt, xrightmargin=0pt]
ON_CALL(precos, preco_de(Eq("caneta"))).WillByDefault(Return(2.50));
ON_CALL(precos, preco_de(_)).WillByDefault(Return(std::nullopt));   // ⚠
\end{lstlisting}
Isto falha com \texttt{C++ exception with description "produto desconhecido:
caneta"}. O gmock tenta as regras da \textbf{mais recente para a mais antiga} ---
e o \texttt{\_}, declarado por último, captura também a caneta. \textbf{A regra
genérica precisa vir primeiro}, para que as específicas, declaradas depois,
ganhem.
\end{armadilha}

Um \emph{mock} tem ainda três níveis de rigor diante de uma chamada que nenhuma
expectativa previu: \texttt{NaggyMock} (o padrão) imprime aviso no \emph{log} e o
teste \emph{passa}; \texttt{NiceMock} ignora em silêncio; \texttt{StrictMock}
\textbf{falha} o teste. Use \texttt{NiceMock} por padrão --- o \emph{naggy} enche
o \emph{log} de ruído e o \emph{strict} quebra o teste toda vez que a
implementação ganha uma chamada legítima a mais.

\begin{regra}[mais de dois ou três \texttt{EXPECT\_CALL} é sinal de teste frágil]
Um teste assim verifica a \emph{implementação}, não o comportamento; ele passará
a falhar em toda refatoração correta. \textbf{Mocke} o que é lento, não
determinístico, irreversível, difícil de levar a um estado de erro, ou o que
ainda não existe. \textbf{Não mocke} tipos de valor, a biblioteca padrão, código
seu que é rápido e determinístico, nem tipos que você não possui --- para esses,
mocke a \emph{sua} camada de adaptação em volta deles.
\end{regra}

\subsubsection{\emph{Death tests}: verificar que o programa morre}

A pré-condição \texttt{assert(quantidade > 0)} aborta o processo quando violada.
Como testar isso, se abortar mataria o processo de teste? O GoogleTest resolve
com \texttt{fork()}: a expressão roda num processo \textbf{filho}, e o pai
verifica o código de saída e a saída em \emph{stderr}.

\begin{lstlisting}[style=bt, language=C++,
    caption={\texttt{EXPECT\_DEATH}, em \texttt{src/testing/test\_death.cpp}.}]
// Uma suíte terminada em `DeathTest` é executada ANTES de todas as outras,
// para reduzir a chance de o fork acontecer com threads já criadas.
TEST_F(CarrinhoDeathTest, QuantidadeZeroAborta) {
    loja::Carrinho c{precos};
    // O segundo argumento casa contra a stderr do filho: é o que distingue
    // "morreu pelo motivo certo" de "morreu por qualquer motivo".
    EXPECT_DEATH(c.adicionar("caneta", 0), "quantidade precisa ser positiva");
}
\end{lstlisting}

\begin{armadilha}[os quatro perigos dos \emph{death tests}]
\textbf{(1) Threads.} \texttt{fork()} num processo com \emph{threads} clona
apenas a \emph{thread} que chamou; se outra segurava um \emph{mutex}, o filho
trava para sempre --- daí a convenção do sufixo \texttt{DeathTest}.
\textbf{(2) Efeitos colaterais.} O que a instrução faz acontece no \emph{filho} e
é descartado: depois de \texttt{EXPECT\_DEATH(\{ x = 1; morre(); \}, "")}, o
\texttt{x} do pai continua zero.
\textbf{(3) O regex.} É POSIX estendido no Linux e um subconjunto simples no
Windows; um regex vazio casa qualquer coisa --- e aí o teste só verifica que
morreu, não por quê.
\textbf{(4) Velocidade.} Cada caso é um \emph{fork} mais uma execução.
\end{armadilha}

Como \texttt{assert} vira um \emph{no-op} sob \texttt{NDEBUG}, os \emph{death
tests} só existem em compilação de depuração; em \emph{release} eles se
auto-pulam com \texttt{GTEST\_SKIP()} --- testar um \texttt{assert} desativado
seria testar o compilador. A macro \texttt{EXPECT\_DEBUG\_DEATH} faz isso
sozinha, dispensando o \texttt{\#ifndef} escrito à mão. E uma exceção
\textbf{não} é morte: é comportamento normal e previsto, para o qual existe
\texttt{EXPECT\_THROW}.

\subsubsection{Cobertura, e o que 100\% não garante}

\textbf{Cobertura} é a fração do código \emph{executada} durante a suíte. Vem em
quatro sabores, e a diferença entre eles é a diferença entre uma métrica
confortável e uma métrica útil.

\begin{table}[htbp]
\centering
\caption{Os quatro tipos de cobertura, para \texttt{if (a \&\& b) faz\_algo();}.}
\label{tab:cobertura-tipos}
\footnotesize
\begin{tabular}{@{}lll@{}}
\toprule
\textbf{Tipo} & \textbf{O que mede} & \textbf{Testes necessários} \\
\midrule
linha    & a linha executou?                      & 1 \\
ramo     & os dois lados do \texttt{if}?          & 2 \\
condição & cada subcondição foi \texttt{true} \textbf{e} \texttt{false}? & \texttt{a} e \texttt{b} variando \\
caminho  & todas as combinações de todos os ramos & exponencial \\
\bottomrule
\end{tabular}
\end{table}

O curto-circuito do \texttt{\&\&} é o que separa \emph{linha} de
\emph{condição}: com \texttt{a} falso, \texttt{b} nunca é avaliado --- a linha
conta como coberta e metade da lógica nunca rodou. A maioria das ferramentas
reporta \emph{linha} por padrão, a métrica mais fácil e menos informativa; peça
ramo com \texttt{gcov -b} ou \texttt{lcov --rc branch\_coverage=1}.

\begin{lstlisting}[style=bt, language=bash,
    caption={Medindo cobertura com GCC e com o Meson.}]
g++ --coverage -O0 -g fonte.cpp teste.cpp -lgtest_main -lgtest
./a.out && gcov -b fonte.cpp        # relatório por arquivo, com ramos
lcov --capture -d . -o cov.info && genhtml cov.info -o html/

meson setup build -Db_coverage=true && meson test -C build
ninja -C build coverage-html
\end{lstlisting}

Compile com \texttt{-O0}: com otimização, \emph{inlining} e eliminação de código
morto embaralham a numeração e o relatório passa a mentir nas duas direções.

\begin{armadilha}[o que 100\% de cobertura NÃO garante]
\textbf{(a) Que as asserções existem.} \texttt{TEST(X, Cobre) \{ minha\_funcao(1,
2); \}} dá 100\% da função e não verifica \textbf{nada}.
\textbf{(b) Que os valores certos foram testados.} Cobrir \texttt{if (n >= 10)}
com \texttt{n=5} e \texttt{n=50} é 100\% de ramo --- e não descobre que o certo
era \texttt{> 10}. As \textbf{fronteiras} encontram esse bug, e cobertura não
pede por elas.
\textbf{(c) Que o código que falta existe.} Um \texttt{else} esquecido não
aparece em relatório nenhum: cobertura mede o que está escrito, não o que deveria
estar.
\textbf{(d) Ausência de \gls{ub}.} Um \texttt{int} que estoura executa a linha e
conta como coberta.
\end{armadilha}

\begin{regra}[cobertura é detector de buracos, não meta numérica]
Lida como relatório --- ``este \texttt{catch} nunca rodou'', ``esta função
inteira está em zero'' --- é uma das ferramentas mais úteis que existem. Como
\emph{gate} numérico do CI, produz exatamente o tipo de teste que o item (a)
descreve. A métrica melhor é o \textbf{teste de mutação}: a ferramenta altera o
código (troca \texttt{>=} por \texttt{>}, inverte um \texttt{if}) e verifica se
\emph{algum} teste falha --- um mutante que sobrevive é um buraco real. Em C++:
\texttt{mull}, \texttt{Dextool}.
\end{regra}

Além dos testes de exemplo, que verificam \emph{um} caso que você imaginou, há
duas técnicas para os casos que você não imaginou. O \textbf{teste baseado em
propriedades} troca ``\texttt{f(2,3) == 5}'' por uma propriedade que vale
\emph{sempre}, e deixa a ferramenta gerar centenas de entradas: idempotência
(\texttt{ordena(ordena(v)) == ordena(v)}), ida e volta
(\texttt{desserializa(serializa(x)) == x}), invariante ou oráculo. No carrinho:
\emph{``o total nunca é maior que o subtotal''}. Biblioteca: RapidCheck. O
\textbf{fuzzing} alimenta a função com entradas geradas e guiadas por cobertura,
procurando \emph{crash} e UB --- vale para tudo que recebe entrada \textbf{não
confiável}: analisadores sintáticos, decodificadores, protocolos.

\begin{lstlisting}[style=bt, language=C++,
    caption={O ponto de entrada do libFuzzer.}]
extern "C" int LLVMFuzzerTestOneInput(const uint8_t* d, size_t n) {
    meu_parser(std::string_view{reinterpret_cast<const char*>(d), n});
    return 0;
}   // clang++ -fsanitize=fuzzer,address fuzz.cpp
\end{lstlisting}

E a combinação de melhor retorno em C++ é rodar a suíte que já existe sob
\emph{sanitizer}: os testes já exercitam o código, e sob instrumentação cada um
vira um detector de \emph{use-after-free}, \emph{overflow}, vazamento e UB ---
\textbf{sem escrever um teste a mais}. É a ponte para a próxima seção.

% =============================================================================
\subsection{Comportamento indefinido}
\label{sec:ub}
% =============================================================================

\begin{porque}[que problema esta seção resolve]
Uma crença: a de que um programa com \gls{ub} ``dá um valor errado''. Não dá. Ele
autoriza o compilador a \textbf{remover código que você escreveu} --- inclusive a
verificação que você pôs justamente para se proteger. A diferença entre as duas
visões é a diferença entre um bug e uma classe inteira de vulnerabilidades.
\end{porque}

A definição do padrão é curta e implacável: \emph{``behavior for which this
document imposes no requirements''}. Nenhuma exigência. E --- o ponto que quase
todo mundo perde --- o efeito \textbf{não precisa se limitar ao trecho errado}:
pode aparecer \emph{antes} da linha do erro, porque o compilador reorganizou o
código supondo que o UB não aconteceria. O exemplo canônico não é acadêmico:

\begin{lstlisting}[style=bt, language=C++,
    caption={O \texttt{tun\_chr\_poll} do kernel Linux --- CVE-2009-1897.}]
struct sock* sk = tun->sk;      // ← desreferencia tun
if (!tun) return POLLERR;       // ← checa DEPOIS
\end{lstlisting}

O GCC raciocinou: \texttt{tun->sk} só é legal se \texttt{tun != nullptr}, logo
\texttt{!tun} é sempre falso, logo o \texttt{if} inteiro pode sumir. E sumiu.
Virou escalação de privilégio --- um atacante conseguia mapear a página zero.

\subsubsection{As quatro categorias, e por que UB existe}

O padrão distingue quatro graus de compromisso, e confundi-los é a origem de
metade das discussões improdutivas sobre portabilidade.

\begin{table}[htbp]
\centering
\caption{As quatro categorias. As três primeiras são seguras; a quarta não tem limite.}
\label{tab:ub-categorias}
\footnotesize
\begin{tabular}{@{}lll@{}}
\toprule
\textbf{Categoria} & \textbf{O compilador\ldots} & \textbf{Exemplo} \\
\midrule
bem definido              & faz o que o padrão manda        & \texttt{2 + 2} \\
\emph{implementation-defined} & escolhe \textbf{e documenta} & \texttt{sizeof(int)} \\
\emph{unspecified}        & escolhe, sem documentar         & ordem de avaliação \\
\textbf{UNDEFINED}        & \textbf{não tem obrigação nenhuma} & \emph{overflow} de \texttt{int} \\
\bottomrule
\end{tabular}
\end{table}

Neste \emph{toolchain}, \texttt{sizeof(long)} é 8 --- e é \textbf{4 no Windows};
\texttt{char} é \texttt{signed} --- e é \texttt{unsigned} em ARM por padrão. Nos
dois casos, o código que supõe o contrário compila e roda até ser portado. A
ordem de avaliação de argumentos é \emph{unspecified}: numa chamada
\texttt{print("\{\} \{\} \{\}", f(), f(), f())} o GCC imprime \texttt{3 2 1}. É
legal, e é armadilha.

\begin{porque}[a troca que o comitê fez de propósito]
\textbf{(a) Desempenho.} Se \texttt{i + 1 > i} pudesse ser falso, o compilador não
poderia simplificar \texttt{for (int i = 0; i <= n; ++i)} nem promover \texttt{i}
a um registrador de 64 bits. Como \emph{overflow} com sinal é UB, ele \emph{pode}
supor que não acontece.
\textbf{(b) Portabilidade.} O padrão precisa valer para máquinas com complemento
de um, palavras de 9 bits e sem unidade de gerenciamento de memória.
\textbf{(c) Custo zero para quem não usa.} Definir ``acesso fora do vetor lança
exceção'' exigiria verificação em \textbf{todo} acesso --- e o C++ não cobra por
segurança que você não pediu.
\end{porque}

O C++20 converteu um caso famoso: inteiros com sinal passaram a ser
obrigatoriamente complemento de dois (P0907). O \emph{overflow} deles, no
entanto, \textbf{continua sendo UB} --- a representação foi fixada, o
\emph{overflow} não.

\begin{table}[htbp]
\centering
\caption{O catálogo de UB, agrupado por quem tem chance de pegá-lo.}
\label{tab:ub-catalogo}
\footnotesize
\begin{tabular}{@{}ll@{}}
\toprule
\textbf{Grupo} & \textbf{Casos} \\
\midrule
memória (ASan) & fora dos limites; \emph{use-after-free}; \emph{use-after-scope};
                 \emph{double free}; \texttt{delete[]} em ponteiro de \texttt{new};
                 \texttt{nullptr}; referência a objeto morto \\
aritmética (UBSan) & \emph{overflow} \textbf{com sinal} (sem sinal é definido!);
                 divisão por zero; \texttt{INT\_MIN / -1}; \emph{shift} $\geq$
                 largura do tipo; \texttt{double} $\rightarrow$ \texttt{int} fora
                 de faixa \\
tipos          & \emph{strict aliasing}; ler membro inativo de \texttt{union};
                 desalinhamento; usar objeto antes do construtor ou depois do
                 destrutor \\
controle       & função não-\texttt{void} sem \texttt{return}; \gls{datarace};
                 laço infinito sem efeito observável; exceção escapando de
                 \texttt{noexcept} \\
biblioteca     & \texttt{v[i]} com \texttt{i >= size()}; desreferenciar
                 \texttt{end()} ou iterador inválido; \texttt{.front()} em
                 container vazio; \texttt{memcpy} com \texttt{nullptr}, mesmo com
                 tamanho 0 \\
\bottomrule
\end{tabular}
\end{table}

Dois itens surpreendem sempre: \emph{overflow} de inteiro \textbf{sem} sinal é
\textbf{bem definido} --- envolve módulo $2^N$, por exigência do padrão; é só o
com sinal que é UB. E ler um membro inativo de uma \texttt{union} é o idioma
padrão de \emph{type punning} em C desde o C99, e é \textbf{UB em C++}.

\subsubsection{Como o UB licencia o otimizador}

A cadeia de raciocínio do compilador é impecável e sempre a mesma: \emph{``este
programa não tem UB --- o padrão diz que programas corretos não têm. Logo, se
\texttt{p} foi desreferenciado, \texttt{p != nullptr}. Logo, a verificação
\texttt{if (p)} três linhas abaixo é sempre verdadeira. Logo, posso apagá-la.''}

\begin{armadilha}[a verificação de \emph{overflow} que não verifica]
\begin{lstlisting}[style=bt, language=C++, frame=none,
                   xleftmargin=0pt, xrightmargin=0pt]
if (a + b < a) trata_overflow();     // ✗ some com -O2
\end{lstlisting}
Se \texttt{a + b} estoura, o programa já tem UB antes da comparação; logo o
compilador pode supor que não estoura; logo \texttt{a + b < a} é sempre falso;
logo a chamada é código morto. As formas corretas verificam \textbf{antes} de
estourar:
\begin{lstlisting}[style=bt, language=C++, frame=none,
                   xleftmargin=0pt, xrightmargin=0pt]
if (__builtin_add_overflow(a, b, &r)) trata_overflow();  // usa o flag da CPU
if (b > 0 && a > INT_MAX - b)         trata_overflow();  // sem extensão
\end{lstlisting}
\end{armadilha}

O arquivo \texttt{src/undefined\_behavior/ub\_exemplos.cpp} --- o único do
repositório que contém UB de propósito, não instalado e com cada caso exigindo um
argumento explícito --- tem doze casos. Um deles produziu o resultado que mais
ensina de todo o subprojeto:

\begin{lstlisting}[style=bt, language=C++,
    caption={O caso \texttt{uaf}: leitura de memória já liberada.}]
int caso_uaf(int n) {
    int* p = new int[4];
    p[0] = n;
    delete[] p;
    return p[0];            // ← use-after-free
}
\end{lstlisting}

\begin{lstlisting}[style=bt, language=bash,
    caption={Medido em \texttt{src/undefined\_behavior/} com
             \texttt{-fsanitize=address}: o MESMO código, dois níveis.}]
-O0 → ERROR: AddressSanitizer: heap-use-after-free
-O1 → imprime 42, sem diagnóstico nenhum
\end{lstlisting}

Com \texttt{-O1}, o compilador viu \texttt{p[0] = n; ... return p[0]} e propagou
a constante: \texttt{return n}. \textbf{A leitura desapareceu antes de o
\emph{sanitizer} poder instrumentá-la.} Não há acesso, logo não há o que
reportar.

\begin{regra}[duas conclusões que valem para qualquer projeto]
(1)~Rode a suíte sob \emph{sanitizer} em \texttt{-O0} \textbf{e} em
\texttt{-O1}/\texttt{-O2}: cada nível esconde e revela coisas diferentes.
(2)~\textbf{``O ASan não acusou'' não é prova de que não há UB.} A recomendação
oficial do ASan é \texttt{-O1}, pelo equilíbrio entre velocidade e fidelidade ---
e este exemplo mostra que ela também tem custo.
\end{regra}

O padrão não diz que a execução fica indefinida \emph{a partir} do UB: diz que a
execução \textbf{inteira} fica indefinida. É o chamado \emph{time travel}:

\begin{lstlisting}[style=bt, language=C++,
    caption={O efeito que aparece antes da causa.}]
void f(int* p, bool cond) {
    printf("antes\n");      // ← pode NÃO ser impresso
    if (cond) *p = 1;       // ← o UB está aqui
}
\end{lstlisting}

Se o compilador provar que \texttt{cond} é sempre verdadeiro e \texttt{p} é
sempre nulo, o corpo inteiro vira ``inalcançável'' e ele pode remover tudo ---
inclusive o \texttt{printf} anterior. É o que torna a depuração de UB tão
desagradável: \textbf{a última mensagem que apareceu não delimita onde o problema
está}. O corolário é que \texttt{std::unreachable()} (C++23) e
\texttt{[[assume(x)]]} são \textbf{UB deliberado}: você promete que aquilo não
acontece, e o compilador otimiza em cima da promessa.

\begin{armadilha}[três casos famosos, e o aviso que pega cada um]
\textbf{\texttt{memset} para apagar senha} --- não é UB, é eliminação de escrita
morta perfeitamente legítima: ninguém lê o \emph{buffer} depois, logo a escrita é
removida, logo a senha fica na memória. Use \texttt{explicit\_bzero} ou um
ponteiro \texttt{volatile}.
\textbf{\texttt{if (this == nullptr)}} --- chamar um método já pressupõe
\texttt{this != nullptr}, então a verificação some
(\texttt{-Wnonnull-compare}).
\textbf{Função sem \texttt{return}} --- cair fora de uma função não-\texttt{void}
é UB, e na prática o compilador emite código que ``continua'' para a função
seguinte no binário. Ligue \texttt{-Werror=return-type}.
\end{armadilha}

\subsubsection{Os \emph{sanitizers}, medidos}
\label{sec:ub-sanitizers}

Um \textbf{\emph{sanitizer}} é instrumentação inserida em tempo de compilação: o
compilador põe verificações em torno de cada acesso, e uma biblioteca de tempo de
execução mantém o mapa do que é memória válida. Custa desempenho e memória, e por
isso não vai para produção --- vai para os \textbf{testes}.

\begin{table}[htbp]
\centering
\caption{Os quatro \emph{sanitizers} e o custo de cada um.}
\label{tab:sanitizers}
\footnotesize
\begin{tabular}{@{}llll@{}}
\toprule
\textbf{Nome} & \textbf{Flag} & \textbf{Pega} & \textbf{Custo} \\
\midrule
Address   & \texttt{-fsanitize=address}   & memória (UAF, \emph{overflow}, vazamento) & ${\approx}2\times$ \\
Undefined & \texttt{-fsanitize=undefined} & UB aritmético e de conversão & ${\approx}1{,}2\times$ \\
Thread    & \texttt{-fsanitize=thread}    & \gls{datarace}   & $5$--$15\times$ \\
Memory    & \texttt{-fsanitize=memory}    & leitura de não inicializado & ${\approx}3\times$ \\
\bottomrule
\end{tabular}
\end{table}

O ASan traz junto o LeakSanitizer, ligado por padrão no Linux: ele reporta, ao fim
do programa, o que não foi liberado.

\begin{armadilha}[ASan e TSan são incompatíveis; MSan é quase inviável]
Os dois primeiros instrumentam os mesmos acessos e não podem ser ligados juntos
--- exigem \textbf{dois jobs separados no CI}. Já o MSan exige que \textbf{toda} a
pilha esteja instrumentada, inclusive a libstdc++: sem uma libc++ compilada com
MSan, ele produz falsos positivos ou --- como medido aqui --- \textbf{simplesmente
não reporta}. É o menos usável dos quatro, e o único que o GCC não implementa. A
combinação padrão, barata e compatível, é \texttt{address,undefined}.
\end{armadilha}

Nenhuma linha da \cref{tab:matriz-ub} foi copiada de documentação: cada uma saiu
de rodar \texttt{ub\_exemplos.cpp} com \texttt{-fsanitize=address,undefined} neste
\emph{toolchain} (GCC 13.3 e clang 18.1.3).

\begin{table}[htbp]
\centering
\caption{Matriz de detecção medida. Três resultados contrariam o senso comum.}
\label{tab:matriz-ub}
\footnotesize
\begin{tabular}{@{}lll@{}}
\toprule
\textbf{Caso} & \texttt{-O0} & \texttt{-O1} \\
\midrule
heap          & ASan: \emph{buffer-overflow}      & ASan: \emph{buffer-overflow} \\
pilha         & UBSan: \emph{index out of bounds} & idem \\
\textbf{uaf}  & \textbf{ASan: \emph{heap-use-after-free}} & \textbf{NADA} \\
dangling      & UBSan: \emph{null reference binding} & idem \\
overflow      & UBSan: \emph{signed integer overflow} & idem \\
shift         & UBSan: \emph{shift exponent too large} & idem \\
divisao       & UBSan: \emph{division by zero}    & idem \\
nulo          & UBSan: \emph{load of null pointer} & idem \\
vazamento     & LSan: \emph{detected memory leaks} & idem \\
\textbf{aliasing}    & \textbf{NADA}              & \textbf{NADA} \\
\textbf{naoinit}     & \textbf{NADA}              & \textbf{NADA} \\
\textbf{alinhamento} & \textbf{NADA (GCC)}        & \textbf{NADA (GCC)} \\
\bottomrule
\end{tabular}
\end{table}

Os três destacados: \textbf{(1)}~\texttt{uaf} some com \texttt{-O1}, pelo motivo
já explicado; \textbf{(2)}~\texttt{alinhamento} não é pego pelo GCC, nem com
\texttt{-fsanitize=alignment} isolado --- o clang pega (\texttt{runtime error:
load of misaligned address ... which requires 4 byte alignment}), de modo que, se
alinhamento importa no seu código, esta verificação exige clang;
\textbf{(3)}~\texttt{aliasing} não é pego por nenhum dos dois.

\begin{armadilha}[o que nenhum \emph{sanitizer} pega]
\textbf{(a) \emph{Strict aliasing}} --- não há instrumentação para isso; a
violação simplesmente produz o valor que o otimizador decidiu.
\textbf{(b) Leitura de não inicializado}, na prática --- o MSan existe e é
inviável sem toda a stdlib instrumentada; a defesa real é
\texttt{-Wmaybe-uninitialized} e a regra de sempre inicializar.
\textbf{(c) Violação de ODR entre unidades de tradução} --- só \texttt{-flto} com
\texttt{-Wodr} tem chance.
\textbf{(d) A lógica errada} --- nenhum \emph{sanitizer} sabe que a sua fórmula de
desconto está trocada; é para isso que existem os testes.
\textbf{(e) O caminho que não foi executado} --- \emph{sanitizer} é análise
\textbf{dinâmica}: só vê o que rodou. Um \texttt{catch} que os testes nunca
exercitam pode ter UB e passar limpo. É por isso que \emph{sanitizer} e cobertura
formam um par: um mede o que rodou, o outro diz o que não rodou.
\end{armadilha}

\begin{lstlisting}[style=bt, language=bash,
    caption={Como rodar --- e a flag que decide se o CI serve para alguma coisa.}]
g++ -std=c++20 -O1 -g -fsanitize=address,undefined \
    -fno-omit-frame-pointer -fno-sanitize-recover=all arq.cpp -o prog

meson setup build-asan -Db_sanitize=address,undefined -Db_lundef=false
meson test -C build-asan

ASAN_OPTIONS=detect_leaks=1:abort_on_error=1
ASAN_OPTIONS=detect_stack_use_after_return=1
UBSAN_OPTIONS=print_stacktrace=1:halt_on_error=1
\end{lstlisting}

\begin{regra}[\texttt{-fno-sanitize-recover=all} é essencial no CI]
Sem ela, o UBSan apenas \textbf{imprime} o diagnóstico e continua a execução --- e
o código de saída do processo fica 0. A suíte passa \emph{com o UB detectado}, e o
relatório verde esconde a mensagem no meio do \emph{log}.
\end{regra}

\subsubsection{\emph{Strict aliasing} e \texttt{bit\_cast}}

A regra, em uma frase: acessar o valor de um objeto por um \emph{lvalue} de tipo
\textbf{incompatível} é UB. As exceções são poucas e vale decorá-las: o tipo exato
(com ou sem \texttt{const}/\texttt{volatile}); a versão com ou sem sinal do mesmo
tipo; uma classe \textbf{base}; e \texttt{char}, \texttt{unsigned char} e
\classe{std::byte} --- os três únicos universais. \texttt{signed char}
\textbf{não} está na lista, e é uma assimetria fácil de esquecer.

\begin{porque}[por que a regra existe]
Para o otimizador. Considere \texttt{void f(int* pi, float* pf)} que escreve em
\texttt{*pi}, escreve em \texttt{*pf} e depois lê \texttt{*pi}. Com \emph{strict
aliasing}, o compilador \textbf{sabe} que \texttt{pf} não pode apontar para o
mesmo objeto que \texttt{pi} e pode reusar o valor que já está no registrador.
Sem a regra, toda escrita por ponteiro invalidaria todos os valores em
registrador, e código com ponteiros ficaria sistematicamente mais lento.
\end{porque}

\begin{lstlisting}[style=bt, language=C++,
    caption={\emph{Type punning}: o errado e os dois certos.}]
int i = *reinterpret_cast<const int*>(&f);      // ✗ UB — e NENHUM sanitizer avisa
std::memcpy(&i, &f, sizeof f);                  // ✓ legal desde sempre
auto i = std::bit_cast<std::uint32_t>(f);       // ✓ C++20

static_assert(std::bit_cast<std::uint32_t>(1.0f) == 0x3F800000u);  // constexpr!
\end{lstlisting}

Medido: os dois caminhos legais dão \texttt{0x4048F5C3} para \texttt{3.14f}. E
\texttt{memcpy} \textbf{não é lento}: com \texttt{-O2} o compilador reconhece o
padrão e emite um único \texttt{mov}. \classe{std::bit\_cast} ainda é melhor em
quatro pontos: exige \texttt{sizeof} igual (erro de \textbf{compilação} se não
for); exige os dois tipos \emph{trivially copyable}; não menciona ponteiro, então
não há como errar o \texttt{sizeof}; e é \texttt{constexpr}.

\begin{armadilha}[\texttt{union} para \emph{type punning}: legal em C, UB em C++]
Ler o membro \textbf{inativo} de uma \texttt{union} é o idioma padrão em C desde o
C99. Em C++ é UB --- só um membro \emph{existe} por vez, e ler outro viola a regra
de tempo de vida. GCC e Clang aceitam como extensão; o MSVC não dá a mesma
garantia. \textbf{Use \texttt{bit\_cast}.}
\end{armadilha}

Quando a violação já está no código antigo e auditar é inviável, existe a saída de
emergência \texttt{-fno-strict-aliasing}, que faz o compilador supor que qualquer
ponteiro pode apontar para qualquer coisa --- o kernel Linux compila com ela desde
2003, não porque a regra seja ruim, mas porque auditar milhões de linhas de C
antigo não é possível. A ordem de preferência é: \texttt{bit\_cast}/\texttt{memcpy};
depois \texttt{char}/\texttt{std::byte} quando você precisa mesmo ver bytes;
depois \texttt{\_\_attribute\_\_((may\_alias))} cirúrgico; e só então a flag no
projeto inteiro.

\begin{armadilha}[a flag precisa valer para TODAS as unidades de tradução]
Uma \textbf{unidade de tradução} é um arquivo \texttt{.cpp} depois do
pré-processador --- a menor coisa que o compilador compila de uma vez. Se uma
delas for compilada \emph{sem} \texttt{-fno-strict-aliasing}, ela otimizará em
cima de uma suposição que outra viola: nenhum erro de \emph{link}, nenhum aviso,
corrupção em execução.
\end{armadilha}

\subsubsection{As cinco camadas de defesa}

Da mais barata para a mais cara --- e a mais barata é também a mais eficaz, o que
costuma surpreender. \textbf{Camada 1: escrever código que não tem como ter UB}
(custo zero), pela substituição sistemática da \cref{tab:camada-um}.

\begin{table}[htbp]
\centering
\caption{Substituições que eliminam categorias inteiras de UB.}
\label{tab:camada-um}
\footnotesize
\begin{tabular}{@{}ll@{}}
\toprule
\textbf{Em vez de\ldots} & \textbf{Use\ldots} \\
\midrule
\texttt{new} / \texttt{delete}        & \classe{std::unique\_ptr}, \classe{std::vector} \\
\texttt{T*} mais tamanho              & \classe{std::span} \\
array de C                            & \classe{std::array} \\
\texttt{v[i]} em código duvidoso      & \texttt{v.at(i)} \\
\texttt{reinterpret\_cast}            & \classe{std::bit\_cast} \\
\texttt{int} para tamanhos            & \texttt{std::size\_t} \\
\texttt{strcpy} / \texttt{sprintf}    & \classe{std::string}, \classe{std::format} \\
\texttt{T*} que pode ser nulo         & \classe{std::optional}, \texttt{T\&} \\
\texttt{union} para \emph{punning}    & \classe{std::variant} / \texttt{bit\_cast} \\
laço com índice                       & \texttt{for (auto\& x : c)} \\
\bottomrule
\end{tabular}
\end{table}

O exemplo que resume a tabela é o \texttt{span}: \texttt{soma\_c(v.data(), 100)}
compila sem reclamar; \texttt{soma\_span(v)} não tem \emph{onde} errar o tamanho.
Sobre \texttt{at()} contra \texttt{[]}: o custo é uma comparação por acesso,
irrelevante fora do laço mais quente.

\textbf{Camada 2: os avisos do compilador} (custo zero). Além de \texttt{-Wall
-Wextra -Wpedantic}, que não são negociáveis: \texttt{-Wshadow},
\texttt{-Wconversion}, \texttt{-Wsign-conversion}, \texttt{-Wold-style-cast},
\texttt{-Wnull-dereference} (o caso do CVE do kernel),
\texttt{-Wnon-virtual-dtor}, \texttt{-Wduplicated-cond}, \texttt{-Wlogical-op}. E
três que devem ser \textbf{erro}: \texttt{-Werror=return-type},
\texttt{-Werror=uninitialized}, \texttt{-Werror=implicit-fallthrough}.

\textbf{Camada 3: análise estática} --- \texttt{g++ -fanalyzer},
\texttt{clang-tidy --checks='bugprone-*,cert-*'}, \texttt{cppcheck
--enable=all}. O valor específico dela é encontrar o caminho que os testes
\textbf{não} executaram, exatamente o ponto cego do \emph{sanitizer}; falso
positivo é a norma, então adote incrementalmente, ligando em código novo.
\textbf{Camada 4: \emph{sanitizers} nos testes} --- a \cref{sec:ub-sanitizers}.
\textbf{Camada 5: \emph{hardening} em produção} ---
\texttt{-D\_GLIBCXX\_ASSERTIONS} (verifica \texttt{v[i]}, \texttt{.front()}),
\texttt{-D\_FORTIFY\_SOURCE=3}, \texttt{-fstack-protector-strong},
\texttt{-fcf-protection}. O primeiro transforma \texttt{v[i]} fora de faixa de UB
em aborto com mensagem, ao custo de uma comparação; Red Hat e Google já o ligam
por padrão.

\begin{armadilha}[\texttt{\_GLIBCXX\_ASSERTIONS} muda a \gls{abi}]
Ele altera o \emph{layout} de alguns tipos da biblioteca padrão. Precisa valer
para \textbf{todas} as unidades de tradução e \textbf{todas} as bibliotecas
linkadas: misturar produz corrupção silenciosa, não erro de \emph{link}. Vale a
mesma disciplina de \texttt{-fno-strict-aliasing}.
\end{armadilha}

\begin{aviso}
A parte honesta: mesmo com as cinco camadas, C++ não é uma linguagem segura por
construção. Elas reduzem a probabilidade; não a zeram. Isso não é argumento para
abandonar C++ --- é argumento para levar as camadas a sério, e para reconhecer que
\textbf{``eu tomo cuidado'' nunca foi uma delas}.
\end{aviso}

% =============================================================================
\subsection{Desempenho: medir, e não supor}
\label{sec:desempenho}
% =============================================================================

\begin{porque}[que problema o subprojeto \texttt{performance} resolve]
A intuição de desempenho que a maioria de nós tem foi formada contando
\emph{operações} --- e a máquina moderna não cobra por operação, cobra por
\textbf{acesso à memória que não estava no cache}. A consequência é que
programadores experientes erram por uma ordem de grandeza ao prever qual de duas
versões é mais rápida. A única saída é medir --- e medir sem se enganar, que é
mais difícil do que parece.
\end{porque}

Este é o subprojeto que mais mede: nada nele é afirmado sem um número obtido na
própria execução, e \textbf{três das medições contrariaram o texto escrito antes
de medir}. Nos três casos o texto estava errado, e a explicação era mais
interessante que a expectativa.

\subsubsection{A hierarquia de memória e a linha de cache}

Um \textbf{cache} é uma memória pequena e rápida, junto ao processador, que guarda
cópias do que foi lido recentemente da RAM. Há três níveis (L1, L2, L3), cada um
maior e mais lento que o anterior.

\begin{figure}[htbp]
\centering
\begin{tikzpicture}[node distance=1.6mm]
  \node[caixaok, minimum width=30mm] (reg) {registrador \quad 0 ciclos};
  \node[caixaok, minimum width=40mm, below=of reg] (l1)  {cache L1 \quad 4};
  \node[caixa,   minimum width=52mm, below=of l1]  (l2)  {cache L2 \quad 12};
  \node[caixa,   minimum width=66mm, below=of l2]  (l3)  {cache L3 \quad 40};
  \node[caixarisco, minimum width=84mm, below=of l3] (ram)
      {RAM \quad 200 ciclos --- $50\times$ mais lento que L1};
\end{tikzpicture}
\caption{A maior assimetria do hardware moderno. Em 200 ciclos o processador
         executaria centenas de instruções --- por isso, em código que percorre
         dados, ``quantas operações'' quase nunca é a pergunta certa, e ``quantos
         \emph{cache miss}'' é.}
\label{fig:hierarquia}
\end{figure}

O processador nunca lê 4 bytes da RAM. Ele lê uma \textbf{linha de cache} --- um
bloco alinhado de 64 bytes no x86-64. Ler \texttt{v[0]} traz \texttt{v[0..15]}
junto, de graça: se você usar os quinze vizinhos, eles saíram grátis; se não usar,
você pagou por eles do mesmo jeito. O experimento soma 1 milhão de elementos de um
vetor de 64\,MB variando o \emph{passo} entre acessos --- mesmo número de somas em
todos os casos:

\begin{lstlisting}[style=bt, language=bash,
    caption={Saída de \texttt{src/performance/}: o joelho está exatamente em
             $64 / \texttt{sizeof(int)}$.}]
passo | bytes entre acessos | tempo
------|---------------------|------------------
    1 |                   4 | 0.38 ms  (1.0×)
    4 |                  16 | 0.36 ms  (0.9×)
    8 |                  32 | 0.61 ms  (1.6×)
   16 |                  64 | 1.39 ms  (3.7×)   ← o joelho
   32 |                 128 | 2.48 ms  (6.5×)
\end{lstlisting}

Até o passo 16, vários acessos compartilham a mesma linha e o custo mal se move; a
partir dali, cada acesso paga uma linha inteira. O passo 32 é ainda pior por causa
do \textbf{prefetcher} --- o circuito que detecta padrões de acesso e busca as
linhas \emph{antes} de serem pedidas: ele traz linhas adjacentes, e com 128 bytes
de passo metade do que ele traz nunca é usada.

\begin{armadilha}[nota de método: um benchmark que produz um número não produziu um número certo]
A primeira versão deste experimento ia até o passo 1024 e afirmava que o tempo
``estabiliza'' depois de 16. A medição desmentiu as duas coisas:
\begin{lstlisting}[style=bt, language=bash, frame=none,
                   xleftmargin=0pt, xrightmargin=0pt]
passo  256 → 1.75 ms   ← CAIU
passo  512 → 1.56 ms   ← caiu mais
passo 1024 → 4.05 ms   ← disparou
\end{lstlisting}
A causa é um defeito do \textbf{experimento}: o \texttt{if (i >= kN) i = 0} faz o
índice dar a volta. Com passo 256, só 65\,536 posições distintas são tocadas ---
4\,MB, que \textbf{cabem no L3}. O benchmark deixou de medir latência de RAM e
passou a medir L3. (O salto em 1024 é outro efeito: 4096 bytes é o tamanho da
página de memória, e ali começa a pressão sobre o TLB, a cache de tradução de
endereços.) A tabela para no passo 32 porque é a faixa em que o experimento mede o
que diz medir.
\end{armadilha}

\subsubsection{Linha contra coluna, \texttt{vector} contra \texttt{list}, e o prefetcher}

Uma matriz $2048 \times 2048$ de \texttt{int} ocupa 16\,MB e é armazenada por
linhas. Somar \emph{todos} os elementos nas duas ordens possíveis:

\begin{lstlisting}[style=bt, language=bash,
    caption={Saída de \texttt{src/performance/}: mesmo $O(n^2)$, mesmo número de
             somas, 35 vezes de diferença.}]
por linha  (i, depois j) :  0.44 ms
por coluna (j, depois i) : 15.54 ms   ← 35× mais lento
\end{lstlisting}

A única diferença é a \textbf{ordem} dos acessos: por coluna, cada acesso salta
8192 bytes --- uma linha de cache nova por elemento, e 15 dos 16 \texttt{int}
trazidos são descartados antes de serem usados. É por isso que BLAS e Eigen fazem
\emph{tiling}: quebram a matriz em blocos que cabem em L1.

\begin{lstlisting}[style=bt, language=bash,
    caption={Saída de \texttt{src/performance/}: somando 2 milhões de \texttt{int}.}]
std::vector<int> : 0.43 ms      4 bytes/elemento, contíguo
std::list<int>   : 2.80 ms     20 bytes/elemento, uma alocação por nó
\end{lstlisting}

A \classe{std::list} paga por elemento: 4 bytes de dado mais dois ponteiros mais o
custo do alocador, cada nó numa alocação separada e potencialmente distante das
outras. O detalhe que fecha o argumento é que ela tem inserção $O(1)$ no meio --- e
ainda assim o \texttt{vector} costuma ganhar \emph{até nesse caso}: o
\texttt{memmove} de $N$ elementos contíguos é mais rápido que \emph{percorrer} a
lista até a posição de inserção. É a palestra de Bjarne Stroustrup em Going Native
2012, e continua valendo (\cref{sec:stl}).

\begin{lstlisting}[style=bt, language=bash,
    caption={Saída de \texttt{src/performance/}: os mesmos 8 milhões de elementos,
             em três ordens.}]
crescente   :  0.97 ms  (1.0×)
decrescente :  1.04 ms  (1.1×)   ← o prefetcher também prevê para trás
aleatória   : 25.69 ms  (26.4×)
\end{lstlisting}

O prefetcher detecta passo constante em qualquer direção. Acesso aleatório não tem
padrão: cada \emph{miss} é pago integralmente. É por isso que uma tabela
\emph{hash} com sondagem linear ganha de uma com encadeamento --- não por fazer
menos operações, mas por o prefetcher conseguir ajudar.

\begin{aviso}
O caso aleatório percorre um \emph{segundo} vetor, o de índices, então parte da
diferença é esse tráfego a mais. O subprojeto registra a ressalva no próprio
código, e ela é o tipo de honestidade que separa uma medição de uma propaganda ---
ainda que o fator seja grande demais para ser só isso.
\end{aviso}

\subsubsection{AoS contra SoA}

\emph{Array of Structs} (AoS) é como quase todo mundo modela, e é a forma mais
legível: um \texttt{vector} de estruturas. \emph{Struct of Arrays} (SoA) inverte:
uma estrutura de vetores, um por campo.

\begin{lstlisting}[style=bt, language=C++,
    caption={Os dois layouts, de \texttt{src/performance/demo\_aos\_soa.hpp}.}]
// AoS — legível, e o padrão
struct Particula { float x, y, z, vx, vy, vz, massa, carga; int id; bool ativa; };
std::vector<Particula> mundo;

// SoA — o layout que a máquina prefere para alguns laços
struct Mundo { std::vector<float> x, y, z, vx, vy, vz, massa, carga; };
\end{lstlisting}

\begin{lstlisting}[style=bt, language=bash,
    caption={Saída de \texttt{src/performance/}: somando APENAS \texttt{massa} de
             4 milhões de partículas.}]
AoS : 3.72 ms      trouxe 152 MB da RAM (a struct inteira, 40 B por partícula)
SoA : 2.47 ms      trouxe  15 MB        (só o vetor de massa)
\end{lstlisting}

Aqui está a primeira lição de método: a razão de \textbf{bytes} é $10\times$, e a
de \textbf{tempo} é só $1{,}5\times$. A diferença entre esses dois números é o
prefetcher --- o acesso do AoS continua sequencial, e o hardware busca as linhas
antes de serem pedidas. Ele não elimina o tráfego; elimina a \textbf{espera} por
ele, até saturar a banda de memória.

\begin{regra}[contar bytes prevê o teto, não o tempo]
O cálculo ``quantos bytes este laço precisa trazer'' dá o limite inferior imposto
pela banda de memória. Quanto do teto se realiza depende do prefetcher, da
vetorização e do número de fluxos simultâneos. Use a conta de bytes para saber
\emph{se vale a pena tentar}; use o cronômetro para saber \emph{o que aconteceu}.
\end{regra}

O demo foi escrito supondo que, num laço que usa \emph{quase todos} os campos, o
AoS venceria --- afinal a \emph{struct} inteira é útil e nada é desperdiçado. Não
venceu:

\begin{lstlisting}[style=bt, language=bash,
    caption={Saída de \texttt{src/performance/}: o laço que usa todos os campos ---
             o resultado que contrariou a hipótese.}]
AoS : 8.62 ms
SoA : 5.39 ms      ← o SoA ainda ganhou (1.60×)
\end{lstlisting}

Duas razões prováveis nesta máquina. Primeira: \texttt{sizeof(Particula)} é
\textbf{40}, e $64 / 40 = 1$ --- cada partícula já ocupa mais de uma linha, então o
AoS não colhe o benefício de ``tudo na mesma linha''. Segunda: o laço SoA
\textbf{vetoriza} (o compilador emite uma instrução SIMD que soma oito
\texttt{float} de uma vez) e o AoS, com passo de 40 bytes, não --- carregar oito
massas espalhadas exigiria um \emph{gather}, que é lento e nem sempre é escolhido.

\begin{regra}[SoA não é ``o layout rápido'']
Ele troca \textbf{desperdício de linha} por \textbf{mais fluxos simultâneos}. O
prefetcher acompanha um número limitado de \emph{streams} --- tipicamente 8 a 16
--- e oito vetores percorridos em paralelo já disputam esse recurso. Para o AoS
ganhar, a \emph{struct} precisaria caber confortavelmente numa linha, o que torna
reordenar campos e cortar \emph{padding} um pré-requisito, e não um detalhe.
\end{regra}

O custo do SoA é legibilidade: some a entidade como objeto --- não existe mais um
\texttt{Particula\&} para passar adiante, e funções que operavam sobre uma
partícula passam a receber um índice e uma referência ao mundo inteiro. Os
meios-termos, em ordem de esforço: \textbf{(1)~separar por temperatura} --- campos
quentes numa \emph{struct}, frios em outra, o que mantém a legibilidade e resolve a
maior parte; \textbf{(2)~AoSoA} --- blocos de $N$ em SoA dentro de um array
(\texttt{struct Bloco \{ float x[8], y[8], z[8]; \}}), que vetoriza bem e mantém
localidade, e é o layout de \emph{engines} de física e de código de alto
desempenho; \textbf{(3)~uma visão} sobre o SoA, um \emph{proxy} que se parece com
\texttt{Particula\&}. E a regra que vale antes de todas: \textbf{só reorganize o
layout depois que o \emph{profiler} apontar aquele laço}.

\subsubsection{\emph{False sharing}: $11{,}7\times$ em código correto}

Duas \emph{threads}, duas variáveis \textbf{diferentes}, nenhum \emph{mutex},
nenhum \gls{datarace}. E o programa fica quase doze vezes mais lento.

\begin{porque}[por que variáveis independentes brigam]
Porque a unidade de coerência de cache não é a variável --- é a \textbf{linha}. Os
núcleos mantêm consistência entre seus caches por um protocolo (MESI) que opera em
linhas inteiras. Se \texttt{a} e \texttt{b} caem na mesma linha de 64 bytes e o
núcleo 0 escreve em \texttt{a}, o protocolo \textbf{invalida a linha} no núcleo 1
--- que precisa buscá-la de novo para ler \texttt{b}. Os dois passam a trocar a
mesma linha entre os caches, ida e volta, milhões de vezes.
\end{porque}

\begin{lstlisting}[style=bt, language=bash,
    caption={Saída de \texttt{src/performance/}: 4 threads × 20 milhões de
             incrementos, cada uma no SEU contador.}]
contadores na MESMA linha : 336.32 ms
contadores SEPARADOS      :  28.72 ms   ← 11.7× mais rápido
\end{lstlisting}

\textbf{O código está correto nos dois casos.} A única diferença entre eles é um
\texttt{alignas(64)}:

\begin{lstlisting}[style=bt, language=C++,
    caption={As duas versões, de \texttt{src/performance/demo\_false\_sharing.hpp}.}]
struct Grudados {                       // 4 × 8 B = 32 B: tudo na mesma linha
    std::atomic<long> c[kThreads]{};
};
struct Separados {
    struct alignas(kLinha) Isolado {    // kLinha = 64
        std::atomic<long> valor{0};
    };
    Isolado c[kThreads]{};
};
\end{lstlisting}

As correções, em ordem: \textbf{(1)~a melhor é não compartilhar} --- acumule numa
variável local da \emph{thread} e escreva uma vez no fim; o \texttt{++local} vira
um registrador, sem tocar em memória nenhuma;
\textbf{(2)~\texttt{alignas(std::hardware\_destructive\_interference\_size)}},
quando o compartilhamento é inevitável --- o custo aqui foi 32 para 256 bytes,
irrelevante para quatro contadores e proibitivo para quatro milhões;
\textbf{(3)~\emph{padding} manual}, se for preciso portabilidade estrita.

\begin{armadilha}[onde o \emph{false sharing} se esconde]
\textbf{(a)} Contadores por \emph{thread} num array: \texttt{std::atomic<long>
acessos[kThreads];}.
\textbf{(b)} Uma \emph{struct} com campos ``independentes'':
\texttt{\{ atomic<size\_t> cabeca; atomic<size\_t> cauda; \}} --- é exatamente por
isso que a fila SPSC de \cref{sec:concorrencia} separa os dois índices com
\texttt{alignas}.
\textbf{(c)} Um \texttt{std::vector<bool>} escrito por várias \emph{threads} ---
aqui é \emph{pior}: os bits compartilham bytes, então é \gls{datarace} de verdade,
UB e não só lentidão.
\textbf{(d)} Objetos pequenos alocados em sequência e distribuídos entre
\emph{threads}: o alocador os põe lado a lado.
Para diagnosticar em vez de adivinhar: \texttt{perf c2c record} e \texttt{perf c2c
report} mostram \emph{quais} linhas estão pingando entre núcleos.
\end{armadilha}

\subsubsection{Predição de desvio, e o experimento que não reproduz mais}

Um processador moderno tem um \emph{pipeline} de 15 a 20 estágios --- começa a
executar uma instrução muitos ciclos antes de terminar a anterior. Ao chegar num
\texttt{if}, não pode esperar a condição: \textbf{aposta} num caminho e continua
especulativamente. Se acerta, custo zero; se erra, o \emph{pipeline} inteiro é
descartado e recomeçado --- 15 a 20 ciclos perdidos.

O experimento clássico soma condicionalmente 8 milhões de valores aleatórios,
ordenados contra embaralhados: no vetor ordenado a condição muda uma única vez e o
preditor acerta quase sempre; no embaralhado ela é aleatória, com $\approx$50\% de
erro.

\begin{lstlisting}[style=bt, language=bash,
    caption={Medido em \texttt{src/performance/}: \texttt{for (int x : v) if (x >=
             128) s += x;} em quatro níveis de otimização.}]
nível | ordenado | embaralhado | razão | cmov no binário
------|----------|-------------|-------|----------------
-O0   | 41.50 ms |   69.31 ms  | 1.7×  |  18
-O1   |  3.40 ms |    3.41 ms  | 1.0×  |  45
-O2   |  3.39 ms |    3.42 ms  | 1.0×  |  36
-O3   |  1.38 ms |    1.38 ms  | 1.0×  |  52
\end{lstlisting}

O efeito aparece em \texttt{-O0} e \textbf{desaparece a partir de \texttt{-O1}}. A
causa está na última coluna: com otimização, o GCC 13 converte o \texttt{if} em
\texttt{cmov} --- \emph{conditional move}, uma instrução que escolhe entre dois
valores \textbf{sem desviar}. Não há aposta, logo não há como errar a aposta. Esta
é a resposta mais votada de todo o Stack Overflow (\emph{``Why is processing a
sorted array faster than an unsorted array?''}, 2012) --- e, no compilador de hoje
com otimização ligada, \textbf{ela não reproduz mais}.

\begin{regra}[um benchmark medido em \texttt{-O0} não diz nada sobre o binário que você entrega]
Duas lições, e a segunda vale para sempre: (1)~o compilador moderno já elimina boa
parte dos desvios imprevisíveis por conta própria; (2)~medir no nível de
otimização de \textbf{produção} não é detalhe de configuração, é a diferença entre
um resultado e um artefato. O efeito de predição continua real --- ele só não está
mais \emph{neste} laço; um desvio que o compilador não consegue converter, com
chamada de função no corpo, continua pagando os ${\approx}15$ ciclos por erro.
\end{regra}

A mesma medição resolve a tentação de reescrever \texttt{if} como aritmética:

\begin{lstlisting}[style=bt, language=bash,
    caption={Saída de \texttt{src/performance/}: as três formas são a mesma coisa,
             porque o compilador já fazia.}]
(a) if (x >= 128) s += x;      3.42 ms
(b) s += x & -(x >= 128);      3.40 ms  (0.99×)
(c) s += x * (x >= 128);       3.35 ms  (0.98×)
\end{lstlisting}

Antes de trocar legibilidade por aritmética, confirme em vez de supor:
\texttt{g++ -O2 -S -masm=intel arq.cpp -o - | grep -c cmov}. As dicas explícitas
--- \texttt{[[likely]]}, \texttt{[[unlikely]]}, \texttt{\_\_builtin\_expect} ---
ajudam quando você \textbf{sabe} a proporção; chutar errado piora, e o \textbf{PGO}
(\texttt{-fprofile-generate} seguido de \texttt{-fprofile-use}) \emph{mede} em vez
de chutar.

\begin{armadilha}[quando não se preocupar com predição de desvio]
Ela só importa quando \textbf{três} coisas coincidem: o desvio está num laço muito
quente; a condição é genuinamente imprevisível; \textbf{e} o corpo é curto --- se
ele faz trabalho de verdade, os ${\approx}15$ ciclos desaparecem no meio. Desvios
que o preditor acerta quase sempre: a condição de parada de um \texttt{for} (erra
\emph{uma} vez, no fim), a verificação de erro que quase nunca dispara, e qualquer
desvio fora de um laço quente. Para medir: \texttt{perf stat -e
branches,branch-misses ./prog} --- abaixo de 1\% é normal, acima de 5\% num laço
quente investigue, acima de 20\% você achou.
\end{armadilha}

\subsubsection{Como medir sem se enganar}

Esta é a parte mais importante do subprojeto: as anteriores dão números, esta
ensina a desconfiar deles. São cinco jeitos de produzir um número falso.

\paragraph{(1) O código foi eliminado.} Se o resultado de um cálculo nunca é lido,
o compilador tem licença para apagar o cálculo inteiro:

\begin{lstlisting}[style=bt, language=bash,
    caption={Saída de \texttt{src/performance/}: o número falso mais comum de
             todos.}]
somando 10 M ints, resultado DESCARTADO :  0.000 ms
o mesmo, com `nao_otimize(s)`           :  1.317 ms
→ o primeiro foi ELIMINADO: ~20000× "mais rápido" porque não aconteceu
\end{lstlisting}

A defesa é um \textbf{sumidouro} --- uma construção que impede o compilador de
provar que o valor é inútil. O bloco de \emph{assembly} abaixo está vazio e não
emite instrução nenhuma; o que ele faz é \emph{informar} ao compilador que uma
função leu \texttt{valor} e pode ter mexido em qualquer memória, e isso basta para
que o cálculo não possa ser removido nem reordenado por cima dele. É o mesmo
truque de \texttt{benchmark::DoNotOptimize} do Google Benchmark.

\begin{lstlisting}[style=bt, language=C++,
    caption={O sumidouro e o cronômetro, de \texttt{src/performance/bench.hpp}.}]
template <typename T>
inline void nao_otimize(T const& valor) {
    asm volatile("" : : "r,m"(valor) : "memory");
}

template <typename F>
double mediana_ms(F&& corpo, int repeticoes = 7) {
    std::vector<double> amostras;
    corpo();                                     // aquecimento, fora da conta
    for (int i = 0; i < repeticoes; ++i) {
        const auto t0 = std::chrono::steady_clock::now();
        corpo();
        const auto t1 = std::chrono::steady_clock::now();
        amostras.push_back(
            std::chrono::duration<double, std::milli>(t1 - t0).count());
    }
    std::sort(amostras.begin(), amostras.end());
    return amostras[amostras.size() / 2];        // MEDIANA, não média
}
\end{lstlisting}

\paragraph{(2) A primeira execução paga a mais.} Falhas de página --- a primeira
escrita numa página recém-alocada exige que o sistema operacional a mapeie de fato
---, cache frio, TLB frio, carregamento de bibliotecas. Sempre \textbf{aqueça}
antes de cronometrar.

\paragraph{(3) Uma amostra é ruído.} Uma interrupção do sistema operacional no
meio da medição distorce a \textbf{média} e não move a \textbf{mediana}. Use a
mediana de 5 a 11 execuções --- e o \classe{std::steady\_clock}, não o
\texttt{system\_clock}, que pode andar para trás quando o NTP ajusta a hora.

\paragraph{(4) O nível de otimização errado} --- é o caso da predição de desvio,
$1{,}7\times$ em \texttt{-O0} e $1{,}0\times$ em \texttt{-O2}.
\textbf{(5) O dado é amigável demais} --- é o caso do passo 256 que cabia no L3.

\begin{regra}[\emph{benchmark} ou \emph{profiler}? Os dois, nesta ordem]
O \emph{benchmark} responde ``a versão B é mais rápida que a A?''; o
\emph{profiler} responde ``\textbf{onde} o programa gasta o tempo?''.
\emph{Profiler} primeiro, para achar o trecho que importa; \emph{benchmark}
depois, para comparar alternativas naquele trecho. Otimizar sem \emph{profiler} é
o erro mais caro que existe --- não porque a otimização não funcione, mas porque
ela se aplica ao lugar errado: acelerar em $10\times$ uma função que consome 2\%
do tempo melhora o total em \textbf{1,8\%}.
\end{regra}

\begin{lstlisting}[style=bt, language=bash,
    caption={\texttt{perf}, o mínimo que vale saber.}]
perf stat ./prog                 # IPC, branch-misses, cache-misses
perf record -g ./prog            # -g = pilha de chamadas
perf report                      # navegável
perf annotate                    # linha a linha, com o assembly
perf stat -e LLC-load-misses ./prog
perf c2c record ./prog           # false sharing
\end{lstlisting}

O \textbf{IPC} (\emph{instructions per cycle}) é o indicador mais informativo de
uma linha só: abaixo de 1,0 significa que o processador está esperando memória;
acima de 2,0 é bom.

\begin{armadilha}[compile com \texttt{-g} mesmo em \emph{release}]
Sem informação de depuração e sem \texttt{-fno-omit-frame-pointer}, o
\texttt{perf report} mostra endereços em vez de nomes de função. \texttt{-g} não
deixa o binário mais lento: só maior. É o que o \texttt{meson.build} deste
subprojeto faz --- ele é o único do repositório que \textbf{fixa} a otimização,
porque o padrão do Meson vem do \texttt{buildtype} (\emph{debug} implica
\texttt{-O0}) e um \emph{benchmark} sem otimização não diz nada sobre o binário
que se entrega. Em WSL e em contêineres, o \texttt{perf} costuma exigir
\texttt{sudo sysctl kernel.perf\_event\_paranoid=1}, e alguns eventos de hardware
simplesmente não existem em máquina virtual.
\end{armadilha}

\begin{table}[htbp]
\centering
\caption{A ordem certa de otimizar: do que mais rende para o que menos rende.}
\label{tab:ordem-otimizar}
\footnotesize
\begin{tabular}{@{}lll@{}}
\toprule
\textbf{Passo} & \textbf{O que é} & \textbf{Ganho típico} \\
\midrule
1. o algoritmo        & $O(n \log n)$ no lugar de $O(n^2)$ & ilimitado \\
2. fazer menos        & cache de resultado, saída antecipada, lote & alto \\
3. o layout dos dados & é onde estão os $10\times$ e $35\times$ deste capítulo & alto \\
4. paralelizar        & limitado pelos núcleos; nova classe de bugs & médio \\
5. micro-otimização   & SIMD manual, código sem desvio & 10--30\% \\
\bottomrule
\end{tabular}
\end{table}

E as \emph{flags}, que são de graça: \texttt{-O2} (o padrão sensato);
\texttt{-O3} (mais vetorização --- \textbf{meça}, às vezes piora);
\texttt{-march=native} (libera AVX2/AVX-512 da máquina local, e produz binário que
\textbf{não roda} em CPU mais antiga); \texttt{-flto}; e o \textbf{PGO}, que é o
ganho mais consistente (5 a 15\%) e quase nunca usado.

\begin{regra}[meça, mude, meça de novo]
Sem o primeiro passo você não sabe onde mexer; sem o terceiro, não sabe se
melhorou. Todo o resto desta seção é consequência dessas duas frases.
\end{regra}

% =============================================================================
\subsection{Síntese}
% =============================================================================

As três seções tratam de assuntos diferentes e chegam à mesma conclusão, que vale
nomear: \textbf{em C++, a ausência de sintoma não é evidência de correção}.

O teste que passa não prova que o código está certo --- prova que aquele caminho,
com aqueles valores, produziu aquele resultado; e 100\% de cobertura não garante
sequer que existam asserções. O \emph{sanitizer} que não acusa não prova que não
há UB --- o \emph{use-after-free} medido aqui \textbf{desaparece} em \texttt{-O1},
e \emph{strict aliasing} não é pego por instrumentação nenhuma. E o
\emph{benchmark} que produz um número não produziu, com isso, um número certo ---
o experimento do passo de cache mediu L3 achando que media RAM, e o experimento
mais citado do Stack Overflow não reproduz mais no compilador de hoje.

A resposta comum aos três é a mesma disciplina: \textbf{várias camadas
independentes, e nenhuma confiança em uma só}. Testes mais cobertura mais
\emph{sanitizer} mais análise estática, porque cada um vê o que os outros não
veem. Medição em vários níveis de otimização, porque cada nível esconde coisas
diferentes. Mediana em vez de média, aquecimento antes do cronômetro, sumidouro em
volta do resultado.

\begin{regra}[o que levar deste capítulo]
Três frases, uma por seção. \textbf{Código difícil de testar é código mal
desenhado} --- o teste é o primeiro cliente a reclamar do seu projeto, e a injeção
de dependência é o que separa uma classe testável em microssegundos de uma que
exige um banco de dados. \textbf{UB não é um valor errado, é a remoção do código
que você escreveu} --- e a defesa mais barata, escrever código que não tem como ter
UB, é também a mais eficaz. \textbf{O layout dos dados vem antes da instrução} ---
os fatores de $35\times$ e $11{,}7\times$ medidos aqui não vieram de trocar
operação por operação, vieram de mudar a ordem em que a memória é tocada.
\end{regra}
